\documentclass[11pt]{article}

\usepackage[margin=1in]{geometry}
\usepackage{amsmath,amssymb,amsthm,mathtools}
\usepackage{enumitem}
\usepackage{hyperref}

\newtheorem{theorem}{Theorem}
\newtheorem{lemma}{Lemma}
\newtheorem{corollary}{Corollary}
\newtheorem{proposition}{Proposition}
\newtheorem{definition}{Definition}
\newtheorem{remark}{Remark}

\DeclareMathOperator{\OPT}{OPT}
\newcommand{\E}{\mathbb{E}}
\newcommand{\Prb}{\mathbb{P}}
\newcommand{\1}{\mathbf{1}}
\newcommand{\cost}{\operatorname{cost}}
\newcommand{\supp}{\operatorname{supp}}
\newcommand{\polylog}{\operatorname{polylog}}
\newcommand{\dd}{\mathrm{d}}
\newcommand{\optpost}{\OPT_{\mathrm{post}}}

\title{Boosted Backbone, Fractional Circulation, and Polylogarithmic Rounding}
\author{}
\date{}

\begin{document}
\maketitle

\begin{abstract}
This note eventually gives an "integral" rounding of a circulation on a backbone 
sampled on $V$. We hope  this integral circulation can be transformed into an 
adaptive policy with bounded cost.
\end{abstract}

\section{Model and the circulation LP}

Let $(\{r\}\cup V,d)$ be a directed metric, where $r$ is the depot.  Write
$N:=\{r\}\cup V$ and let
\(
  n:=\max\{|N|,3\}.
\)
Each client $v\in V$ is active independently with probability $p_v$.  For an
active set $A\subseteq V$, let $\OPT(A)$ be the length of an optimal directed
TSP tour on $\{r\}\cup A$.  Define
\(
  \OPT_{\mathrm{post}}:=\E_A[\OPT(A)].
\)

A backbone is simply a subset of $V$. We use this concept inspired from Shmoys's
paper, where a backbone is sampled with probability $p_v$ (the same
as the activation probability) and then he constructs detour for every vertex 
outside of the backbone. Obviously, the expected cost of backbone tour is the 
same as $\optpost$. And he managed to prove the detour cost is also 
bounded by a constant multiplicative factor of $\optpost$.

We are hoping that we can have a similar algorithm. We first sample a backbone, 
and then we try to create a detour. However, a detour in asymmetric case is not 
as simple as connecting vertex to its closest neighbour in backbone. That's why 
we need the circulation. A circulation doesn't directly provide a way to connect 
all these outside vertices in a low-cost closed walk, even though the circulation 
itself has low cost. This is the gap we need to solve.

For a backbone set $S\subseteq V$, let
\[
  R:=S\cup\{r\},
  \qquad
  U:=V\setminus S.
\]
So here $R$ is the backbone vertices with root, and $U$ is the outside vertices.

We consider the following directed circulation LP:
% We consider the following directed circulation LP, with an additional degree
% upper bound $D$ on the backbone:
\[
\begin{aligned}
\min \quad & \sum_{u,v\in \{r\}\cup V} d(u,v)x_{uv} \\
\text{s.t.}\quad
& x(\delta^+(w))=x(\delta^-(w)) && \forall w\in \{r\}\cup V, \\
& x(\delta^+(v))=x(\delta^-(v)) = p_v && \forall v\in U, \\
% & x(\delta^+(s))=x(\delta^-(s))\le D && \forall s\in R, \\
& x_{uv}=0 && \forall u,v\in U, \\
& x_{uv}\ge 0 && \forall u,v.
\end{aligned}
% \tag{$\mathrm{LP}_D(S)$}
\tag{$\mathrm{LP}(S)$}
\]
Here $x(\delta^+(w)):=\sum_z x_{wz}$ and
$x(\delta^-(w)):=\sum_z x_{zw}$.

We do not allow edges between two outside vertices, so this circulation gives 
a fractional analogy of a detour for an outside vertex: we go from a backbone 
vertex, visit the outside vertex, and then come back to another backbone vertex.

\section{Circulation has low cost}

In this section we prove the optimal solution to LP has expected low cost, 
and since we are only considering analytic results, so as long as the expectation 
is low, we are guaranteed to have at least one good choice of $S$. Indeed, we 
can repeat the random sampling procedure of $S$ for $O(\log n)$ times and we 
have a low-cost solution with $1-o(1)$ probability. We will omit this technicality 
for now.


The intuition is that, consider any realized active subset $A \subseteq V$, we 
can divide the optimal TSP tour on $A \cup R$ by backbone vertices and connect 
all middle vertices to its two surrounding backbone vertex. So if a segment of 
optimal TSP tour on $A \cup R$ looks like $s_1, u_1, u_2, \dots, u_k, s_2$, where
$s_1, s_2$ are the backbone vertices and $u_i$'s are outside vertices, then 
we create flow $(s_1, u_1, s_2), \dots, (s_1, u_k, s_2)$ with weight $1$. Then 
if we weighted-average over all $A$ we immediately satisfy the flow constraint 
for outside vertex. And we can include auxiliary flow from $s_1$ to $s_2$ to 
satisfy the flow constraint for $S$ vertices. 

Formally speaking, we cut the tour $OPT(A \cup R)$ at the vertices from $R$. 
This gives a list of directed segment $P_j = (s_j, u_{j,1}, \dots, u_{j, k_j}, s_{j+1})$ 
where $s_j \in R$ and $u_{j, i} \in A - R$. 
Define $M(A, R) := \max_j k_j$ to be the maximal number of outside vertices 
between two consecutive backbone vertices under this realization $A$ and backbone $R$. 
We create flow $(s_j, u_{j, 1}, s_{j+1}), \dots, (s_{j}, u_{j,k_j}, s_{j+1})$ with weight $1$.
And we create flow $(s_{j}, s_{j+1})$ with weight $M(A,R) - k_j$.
Denote the flow(edge) weight as $x^{A,R}$ then it is straightforward that 
the cost $\cost(x^{A,R}) \le M(A,R) \cost(T_{A \cup R})$.
We further denote $x^{R} = \E_A [x^{A,R}]$ to be averaged flow. Clearly the flow 
satisfies circulation constraints and flow $p_v$ constraints.

As for the cost of $x^{A,R}$, for each $u_i$, we pay weight 1 for $d(s_1, u_i) + d(u_i, s_2)$,
which can be as large as the total length of $(s_1, u_1, \dots, u_k, s_2)$. 
So to prove the circulation does not cost too much, we need to prove that we can
control the cardinality of vertices between two backbone vertices. 


\begin{remark}
    For a pair $(A, R)$, we have $\cost(x^{A,R}) \le M(A,R) \cost(T_{A \cup R})$.
\end{remark}

The next idea is to bound $M(A, R)$ conditioned on $C := A \cup R$. Given 
any subset $r \in C \subseteq V$ and denote $T_C$ to be the optimal tour for $C$.
Then we prove $\E[M(A,R) | C = A \cup R]$ is bounded, thus $\E_{A,R}[M(A,R)] = \E_C [\E[M(A,R) | C = A \cup R]]$ 
is also bounded.

\begin{lemma}[Longest-run]
    $\Prb[M(A,R) \ge t|C=A\cup R] \le n 2^{-t}$.
\end{lemma}
\begin{proof}
    Condition on $C = A \cup R$, the tour $T_C$ is fixed. For each $v \in C$, 
    the event $v$ is an outside vertex is $\Prb[v \in A - R | v \in A \cup R] = 
    \frac{\Prb[v \in A - R]}{\Prb[v \in A \cup R|} = \frac{p_v(1-p_v)}{1-(1-p_v)^2} \le 1/2$.
    So $\Prb[M(A,R) \ge t | C = A \cup R] \le n \; \text{position and contiguous} \; t \le n 2^{-t}$.
\end{proof}
\begin{corollary}
    $\E[M(A,R)|C=A\cup R] \le O(\log n)$, thus
    $\E[M(A,R)] \le O(\log n)$.
\end{corollary}

\begin{proof}
    Let $L = \log n$ and we have
    \begin{align*}
        \E[M(A,R)|C=A \cup R] &\le L \cdot \Prb[M(A,R) \le L | C = A \cup R]  + \sum_{t=L}^{\infty} t \cdot \Prb[M(A,R) \ge t | C = A \cup R] \\
                              &\le L \cdot 1 + \sum_{t=L}^{\infty} t n 2^{-t} \\ 
                              &\le O(\log n)
    \end{align*}
    Then by total expectation 
    \begin{align*}
        \E[M(A,R)] = \E_{C} [\E[M(A,R) | C = A \cup R]] \le O(\log n)
    \end{align*}
\end{proof}


Finally we can show that the expected cost is low.

\begin{lemma}[Expected cost]
\[
    \E_{R} [\cost(x^R)] \le O(\log n) \OPT_{post}
\]
\end{lemma}
\begin{proof}
    \begin{align*}
        \E_{R} [\cost(x^R)] &= \E_{A,R} [\cost(x^{A,R})] && \text{linearity of cost} \\
                            &\le \E_{A,R} [M(A,R) \OPT(A \cup R)] \\ 
                            &= \E_C[\OPT(C) \E_{A,R}[M(A,R) | C = A \cup R]] \\
                            &\le \E_C[\OPT(C) O(\log n)] \\
                            &= O(\log n) \E_{A,R}[\OPT(A \cup R)] \\ 
                            &\le O(\log n) \cdot 2 \cdot \E_A[\OPT(A)] \\ 
                            &= O(\log n) \OPT_{post}
                            % &\le \E_{A,R}[M(A,R)] \E_{A,R}[\OPT(A \cup R)]
    \end{align*}
\end{proof}

For simplicity, we assume we can simply find the best $R$ such that $\cost(x^R)$ 
is minimized, since we are currently considering analytic problem.

Also we can show that with high probability we can find a solution with bounded 
cost together with bounded degree ($x(\delta^{out}(s)) \le O(\log n)$, where $s \in S$),
we omit the proof for now and we do not need this bounded degree property.

\section{Rounding the circulation}
In the previous section, we have proved that for a good backbone $S$, we can 
have a fractional circulation solution such that the cost is bounded by $O(\log n) \OPT_{post}$.
In this section, we provide a way to \emph{round} this fractional solution. 
Technically, this procedure is not a rounding procedure, because we do not round 
solution to integer, instead, we round solution to either tight upper bound or 
lower bound; also, after rounding, the solution is no longer a strict circulation,
but the circulation constraint (incoming equals outgoing) would not be violated 
very much.

Firstly, consider a slightly different way of representing circulation variables.
Previously, $x_{(u,v)}$ represents the flow of edge $(u,v)$, where $u, v$ can be 
backbone vertices or outside vertices. Now we consider $x_{i,(u,v)}$ to be the flow 
of path $u \rightarrow i \rightarrow v$, where $u, v$ are backbone vertices and $i$
is outside vertex. So we are basically coupling the previous solution. 
Notice that in previous solution, we can also have $x_{(u,v)}$ where $u, v$ are 
both backbone vertices, which cannot be represented as $x_{i, (u,v)}$. So we introduce
$y_{(u,v)}$ to be such variables.

Secondly, we change the constraints a little bit. The outside vertices' flow constraints
stay unchanged, we still have the \emph{outside constraints}:
\begin{align*}
    \sum_{e = (u,v)} x_{i,e} = p_i 
\end{align*}
But for backbone vertices, we divide the backbone $R$ using dyadic intervals, like 
a segment tree style. And for every such dyadic interval, we have \emph{interval constraints}:
\begin{align*}
    x(I^{out}) + y(I^{out}) = x(I^{in}) + y(I^{in})
\end{align*}

Now we consider an iterative rounding procedure. Suppose $F$ to be the set of all 
non-tight $x$ variables, i.e. these $x_{i,e}$ such that $0 < x_{i,e} < p_i$.
We want to prove that during the iterative rounding procedure, we always have that 
$|F|$ is larger than the number of tight constraints, thus we can always find two 
directions to move. Then we move to either direction with certain probability 
keeping the expectation cost unchanged.

Suppose there are $c_1$ number of tight outside constraints and $c_2$ number of
tight interval constraints. 
Unfortunately, it is not true that we always have larger $|F|$ than 
$c_1 + c_2$.
To solve this problem, we need to delete some constraints alongside the procedure.
we only consider the interval constraints when the constraints has at least 
$8 \log n$ non-tight $x$ variables. 
Then we know $c_2 \cdot 8 \log n \le |F| \cdot 2 \log n$ because each non-tight
variable will touch at most $2 \log n$ intevals (we are considering only dyadic ones).
Moreover, $2 c_1 \le |F|$ because if such constraint has one non-tight variable,
it must have at least another non-tight variable.
So it is now clear that $c_1 + c_2 < |F|$ always holds during the procedure, 
which means we can find a direction to move and make at least one non-tight 
variable tight. In the end we will have all $x$ variables being tight.

Since we delete constraints alongside the procedure, the resulting solution is 
no longer a strict circulation. But a constraint is deleted only when it touchs 
at most $8 \log n$ non-tight variables, so the constraint will be off the track
by at most $8 \log n$.

Then for any interval (not necessarily dyadic), it can be splited into $O(\log n)$
dyadic intevals thus the outgoing flow and incoming flow would be off by at most 
$O(\log ^2 n)$.
 
So in summary what does the output look like? Let the backbone $R$ be the chosen backbone equipped
with optimal TSP tour $T = (r = s_0, s_1, \dots, s_{m-1}, s_m = r)$. Further let
$B := \sum_{i=1}^{m-1} d(s_i, s_{i+1})$. After rounding, variables $x_{i, (u,v)}$
assures for every outside vertex $i$, exactly one edge $(u,v), u, v \in R$ is selected.
We also make sure the edges cost is bounded by $O(\log n) \OPT_{post}$. Note that
the edge cost is weighted, each edge selected after rounding cost $p_i \cdot (d(u,i)+d(i,v))$.
The rounding solution also assures for each interval the incoming and outgoing 
imbalance before realization is bounded by $O(\log^2 n)$.

\section{What can we get from circulation and rounding}
We do not directly have an adaptive policy yet. After rounding, 
every outside vertex $i$ has exactly one edge $e$ such that $x_{i,e}=p_i$.
A naive policy would be that, when we stand at a backbone vertex $v$, we always 
try to probe from $v$ to some outside vertex $i$ which has $e = (v, v')$ and 
$x_{i,e}=p_i$. If we cannot find any outside vertex from $v$, then we go alonside
the backbone and continue until we probed all outside vertex. This naive policy 
does not provide good cost upper bound because it relies heavily on backbone edges 
to move from one vertex to another. We do not care about log factor,
we can use each backbone edge poly-log times, but it can happen that poly-log is 
not enough. In circulation lemma, we also paid for these arcs on backbone tour,
i.e. those $u \rightarrow i \rightarrow v$ arcs where $i$ is the outside vertex.
Since ciculation lemma gives bounded cost, we can also use these edges poly-log times.
But notice that, each such arc has non-integral weight. So when we say poly-log times,
we are actually saying we can use it no more than poly-log times the flow weight $x$.
Or in another way, since after rounding each edge has flow weight either $0$ or $p_i$,
we can assume we can use an edge when the corresponding outside vertex is active, thus
in expectation point of view it is equivalent. So in summary, we do not care about 
the edge distance / cost, we essentially care about the congestion! We care about 
how many times a edge in circulation and rounding would be used in adaptive policy.
If we can control this to be upper bounded by poly-log, then we are happy.





\end{document}
